% GENERATED FILE. Edit urdu/pronunciation-reference.md, then run npm run generate:books.

\chapter*{Pronunciation \& Script Reference}
\addcontentsline{toc}{chapter}{Pronunciation \& Script Reference}
\markboth{Pronunciation}{Pronunciation}

A \textbf{reference}, not a chapter. Use it when a lesson names a sound; do not clear it as an alphabet gate. Urdu script is introduced through useful expressions, with only the shapes and sound contrasts needed on that page.

The track shows a learner romanization beside new Urdu. A line above a vowel marks length (\textbf{ā, ī}), and a tilde marks nasalization (\textbf{ā̃, ī̃}). Cover the romanization once the Urdu form is familiar.

\section*{Script, spelling, and type style}

\begin{itemize}
\raggedright
\setlength{\itemsep}{0.35em}
  \item \textbf{Read right to left} (\texttt{rtl}). Start each word at its right edge.
  \item \textbf{Letters usually join} and change shape by position. Some letters, including \textbf{\ur{ا}} and \textbf{\ur{ر}} in current words, break the connection to the following shape on their left.
  \item Urdu is an \textbf{abjad}: some vowels are unwritten and inferred from the learned word (\texttt{short-vowels-unwritten}). Long vowels use letter shapes such as \textbf{\ur{ا}}, \textbf{\ur{و}}, and \textbf{\ur{ی}}; \textbf{\ur{ے}} can carry \emph{e/ai}.
  \item Urdu is traditionally printed in \textbf{Nastaliq}, whose connected words descend in a flowing slope. The book and app use the vendored static Noto Nastaliq Urdu family, including its Urdu OpenType localization and contextual joins. Naskh remains an accessibility fallback when a learner's browser cannot load the course font; it is no longer the normal presentation.
\end{itemize}

\section*{Sound ids used by the starter lessons}

\begin{itemize}
\raggedright
\setlength{\itemsep}{0.35em}
  \item \begin{minipage}[t]{\linewidth}
  \textbf{Sound id:} \texttt{rtl}
  \par \textbf{What to do:} follow the word from its right edge leftward
  \par \textbf{First anchor:} \textbf{\ur{سلام}} \emph{salām}
  \end{minipage}
  \item \begin{minipage}[t]{\linewidth}
  \textbf{Sound id:} \texttt{long-a}
  \par \textbf{What to do:} hold \emph{ā} steady; \textbf{\ur{ا}} carries it in \textbf{\ur{سلام}}
  \par \textbf{First anchor:} \textbf{\ur{سلام}} \emph{salām}
  \end{minipage}
  \item \begin{minipage}[t]{\linewidth}
  \textbf{Sound id:} \texttt{alif-madd}
  \par \textbf{What to do:} read \textbf{\ur{آ}} as word-initial long \emph{ā}
  \par \textbf{First anchor:} \textbf{\ur{آپ}} \emph{āp}
  \end{minipage}
  \item \begin{minipage}[t]{\linewidth}
  \textbf{Sound id:} \texttt{short-vowels-unwritten}
  \par \textbf{What to do:} supply the learned short vowels even when vowel marks are absent
  \par \textbf{First anchor:} \textbf{\ur{شکریہ}} \emph{shukriyā}
  \end{minipage}
  \item \begin{minipage}[t]{\linewidth}
  \textbf{Sound id:} \texttt{long-i}
  \par \textbf{What to do:} hold \textbf{\ur{ی}} as \emph{ī} in this word
  \par \textbf{First anchor:} \textbf{\ur{جی}} \emph{jī}
  \end{minipage}
  \item \begin{minipage}[t]{\linewidth}
  \textbf{Sound id:} \texttt{long-u}
  \par \textbf{What to do:} hold \textbf{\ur{و}} as long \emph{ū} in the learned word
  \par \textbf{First anchor:} \textbf{\ur{ہوں}} \emph{hūṅ}
  \end{minipage}
  \item \begin{minipage}[t]{\linewidth}
  \textbf{Sound id:} \texttt{nasal-vowel}
  \par \textbf{What to do:} let air pass through the nose; final \textbf{\ur{ں}} does not add a full English \emph{n}
  \par \textbf{First anchor:} \textbf{\ur{ہاں}} \emph{hā̃}, \textbf{\ur{نہیں}} \emph{nahī̃}
  \end{minipage}
  \item \begin{minipage}[t]{\linewidth}
  \textbf{Sound id:} \texttt{long-vowels}
  \par \textbf{What to do:} read the word's written long-vowel cues rather than stretching every vowel
  \par \textbf{First anchor:} \textbf{\ur{میرا} \ur{نام}} \emph{merā nām}
  \end{minipage}
  \item \begin{minipage}[t]{\linewidth}
  \textbf{Sound id:} \texttt{hai-diphthong}
  \par \textbf{What to do:} say \textbf{\ur{ہے}} as one \emph{hai} syllable, ending in an \emph{ai} glide
  \par \textbf{First anchor:} \textbf{\ur{ہے}} \emph{hai}
  \end{minipage}
  \item \begin{minipage}[t]{\linewidth}
  \textbf{Sound id:} \texttt{sh}
  \par \textbf{What to do:} read three-dotted \textbf{\ur{ش}} as \emph{sh}
  \par \textbf{First anchor:} \textbf{\ur{شکریہ}} \emph{shukriyā}
  \end{minipage}
  \item \begin{minipage}[t]{\linewidth}
  \textbf{Sound id:} \texttt{kh}
  \par \textbf{What to do:} make \textbf{\ur{خ}} farther back than English \emph{h}
  \par \textbf{First anchor:} \textbf{\ur{خوشی}} \emph{khushī}
  \end{minipage}
  \item \begin{minipage}[t]{\linewidth}
  \textbf{Sound id:} \texttt{consonantal-ye}
  \par \textbf{What to do:} let \textbf{\ur{ی}} supply consonantal \emph{y} before a vowel
  \par \textbf{First anchor:} \textbf{\ur{کیا}} \emph{kyā}
  \end{minipage}
  \item \begin{minipage}[t]{\linewidth}
  \textbf{Sound id:} \texttt{ai-diphthong}
  \par \textbf{What to do:} say \emph{ai} as one glide in the learned question word
  \par \textbf{First anchor:} \textbf{\ur{کیسے}} \emph{kaise}
  \end{minipage}
  \item \begin{minipage}[t]{\linewidth}
  \textbf{Sound id:} \texttt{final-ye}
  \par \textbf{What to do:} distinguish broad final \textbf{\ur{ے}} \emph{e} from final \textbf{\ur{ی}} long \emph{ī}
  \par \textbf{First anchor:} \textbf{\ur{کیسے} / \ur{کیسی}} \emph{kaise / kaisī}
  \end{minipage}
  \item \begin{minipage}[t]{\linewidth}
  \textbf{Sound id:} \texttt{retroflex-aspirated-th}
  \par \textbf{What to do:} curl the tongue slightly back for \textbf{\ur{ٹ}}, then release the aspiration marked by \textbf{\ur{ھ}}
  \par \textbf{First anchor:} \textbf{\ur{ٹھیک}} \emph{ṭhīk}
  \end{minipage}
  \item \begin{minipage}[t]{\linewidth}
  \textbf{Sound id:} \texttt{urdu-question-mark}
  \par \textbf{What to do:} recognize \textbf{\ur{؟}} at the left end of a right-to-left question
  \par \textbf{First anchor:} \textbf{\ur{آپ} \ur{کا} \ur{نام} \ur{کیا} \ur{ہے؟}}
  \end{minipage}
  \item \begin{minipage}[t]{\linewidth}
  \textbf{Sound id:} \texttt{be-vs-pe-dots}
  \par \textbf{What to do:} count the dots under the shared low scoop: one below is \textbf{\ur{ب}} \emph{b}, three below are \textbf{\ur{پ}} \emph{p}
  \par \textbf{First anchor:} \textbf{\ur{بولنا}} \emph{bolnā} against \textbf{\ur{آپ}} \emph{āp}
  \end{minipage}
  \item \begin{minipage}[t]{\linewidth}
  \textbf{Sound id:} \texttt{geminate-nun}
  \par \textbf{What to do:} hold the \emph{n} long when two \textbf{\ur{ن}} letters are written in a row
  \par \textbf{First anchor:} \textbf{\ur{جاننا}} \emph{jānnā}
  \end{minipage}
  \item \begin{minipage}[t]{\linewidth}
  \textbf{Sound id:} \texttt{che-vs-jim-dots}
  \par \textbf{What to do:} count the dots below the \emph{jīm} body: one is \textbf{\ur{ج}} \emph{j}, three are \textbf{\ur{چ}} \emph{ch}
  \par \textbf{First anchor:} \textbf{\ur{سوچنا}} \emph{sochnā} against \textbf{\ur{جانا}} \emph{jānā}
  \end{minipage}
  \item \begin{minipage}[t]{\linewidth}
  \textbf{Sound id:} \texttt{do-chashmi-he}
  \par \textbf{What to do:} give \textbf{\ur{ھ}} no sound of its own; let it aspirate the letter to its right
  \par \textbf{First anchor:} \textbf{\ur{سمجھنا}} \emph{samajhnā}, \textbf{\ur{لکھنا}} \emph{likhnā}
  \end{minipage}
  \item \begin{minipage}[t]{\linewidth}
  \textbf{Sound id:} \texttt{retroflex-flap-rre}
  \par \textbf{What to do:} flick a curled-back tongue once for \textbf{\ur{ڑ}}; do not roll it
  \par \textbf{First anchor:} \textbf{\ur{پڑھنا}} \emph{paṛhnā}
  \end{minipage}
  \item \begin{minipage}[t]{\linewidth}
  \textbf{Sound id:} \texttt{majhul-e}
  \par \textbf{What to do:} read medial \textbf{\ur{ی}} as long \emph{e} where the learned word calls for it
  \par \textbf{First anchor:} \textbf{\ur{لینا}} \emph{lenā}
  \end{minipage}
\end{itemize}

\section*{Shape families already in use}

\begin{itemize}
\raggedright
\setlength{\itemsep}{0.35em}
  \item \textbf{\ur{س}} \emph{s} and \textbf{\ur{ش}} \emph{sh} share a skeleton; three dots distinguish \textbf{\ur{ش}}.
  \item \textbf{\ur{ن}} is consonantal \emph{n}. Final dotless \textbf{\ur{ں}} is \emph{nūn-e ghunna} and marks nasalization in \textbf{\ur{ہاں}} and \textbf{\ur{نہیں}}.
  \item \textbf{\ur{ی}} supplies the long \emph{ī} in \textbf{\ur{جی}} and \textbf{\ur{نہیں}}. Final \textbf{\ur{ے}} participates in \emph{hai} in \textbf{\ur{ہے}}. Learn each role in its word before generalizing.
  \item \textbf{\ur{کیا}} gives \textbf{\ur{ی}} its consonantal \emph{y} job; \textbf{\ur{کیسے} / \ur{کیسی}} then contrast broad final \textbf{\ur{ے}} with long-\emph{ī} final \textbf{\ur{ی}}.
  \item \textbf{\ur{ٹھیک}} introduces the Indic retroflex-plus-aspiration sequence \textbf{\ur{ٹھ}} only when the wellbeing reply needs it.
  \item \textbf{\ur{خدا} \ur{حافظ}} reuses \textbf{\ur{خ}} and long \textbf{\ur{ا}}, adds a clear initial \emph{h} in \textbf{\ur{حافظ}}, and reads final \textbf{\ur{ظ}} as \emph{z}. Urdu conventionally keeps the two words visibly spaced, even though Persian normally joins its local spelling.
  \item \textbf{\ur{بولنا}} adds the only new consonant the verb chapter needs: \textbf{\ur{ب}} \emph{b}, one low scoop with a single dot below. It shares that skeleton with \textbf{\ur{پ}} \emph{p}, already read in \textbf{\ur{آپ}}, which carries three dots below instead. Dot count is the whole difference, and it is load-bearing throughout the alphabet.
  \item \textbf{\ur{جاننا}} writes two \textbf{\ur{ن}} letters in a row and pronounces both, which is the only thing separating it from \textbf{\ur{جانا}} \emph{jānā}. Nothing about the doubling is decorative: \textbf{\ur{جانا}} is “to go” and \textbf{\ur{جاننا}} is “to know.”
  \item \textbf{\ur{چ}} \emph{ch} is the \textbf{\ur{ج}} \emph{j} body with \textbf{three dots below} instead of one — the same dot-count contrast as \textbf{\ur{ب}} against \textbf{\ur{پ}}, on a second skeleton.
  \item \textbf{\ur{ھ}} is \emph{do-chashmī he}, “two-eyed he.” It is never a consonant on its own; it aspirates whatever letter stands to its right. \textbf{\ur{ٹھیک}} used it first, and \textbf{\ur{سمجھنا}}, \textbf{\ur{پڑھنا}}, \textbf{\ur{پوچھنا}} and \textbf{\ur{لکھنا}} use it again in \textbf{\ur{جھ}}, \textbf{\ur{ڑھ}}, \textbf{\ur{چھ}} and \textbf{\ur{کھ}}. Keep it apart from round \textbf{\ur{ہ}} \emph{gol he}, which is a real \emph{h} and is the letter in \textbf{\ur{ہوں}} and \textbf{\ur{ہونا}}.
  \item \textbf{\ur{ڑ}} is \textbf{\ur{ر}} \emph{r} wearing the small retroflex mark that also rides on \textbf{\ur{ٹ}}. One diacritic, one instruction — curl the tongue back — now on two letters.
  \item \textbf{\ur{ی}} takes a \textbf{third} value in \textbf{\ur{لینا}} \emph{lenā}: a long \emph{e}. With consonantal \emph{y} in \textbf{\ur{کیا}} and long \emph{ī} in \textbf{\ur{جی}}, and broad final \textbf{\ur{ے}} beside it, the \emph{ye} family covers \emph{y}, \emph{ī} and \emph{e}; the learned word decides.
\end{itemize}

\section*{The Persian-Arabic and Indo-Aryan layers}

Script does not determine a word's history. \textbf{\ur{سلام}} and \textbf{\ur{شکریہ}} expose the Persian-Arabic literary layer; \textbf{\ur{نہیں}}, \textbf{\ur{میرا}}, \textbf{\ur{نام}}, and \textbf{\ur{ہے}} expose the inherited Indo-Aryan core that Urdu shares with Hindi. Mixed-language practice may use that relationship as a bridge, but the Urdu script remains the assessed form for this track.

\section*{Sources}

\begin{itemize}
\raggedright
\setlength{\itemsep}{0.35em}
  \item \href{https://openbooks.library.northwestern.edu/zerozabar/front-matter/introduction/}{\emph{Zero Zabar}: How the Urdu script works}
  \item \href{https://openbooks.library.northwestern.edu/zerozabar/chapter/the-urdu-alphabet/}{\emph{Zero Zabar}: The Urdu alphabet}
\end{itemize}
